\documentclass[11pt,a4paper]{article}

% ============================================================
% FICHA DE TRABALHO — LEI DA GRAVITAÇÃO UNIVERSAL
% EUFISICA — Prof. José Gonçalves
% ============================================================

% ------------------------------------------------------------
% CODIFICAÇÃO E IDIOMA
% ------------------------------------------------------------

\usepackage[utf8]{inputenc}
\usepackage[T1]{fontenc}
\usepackage[portuguese]{babel}
\usepackage{lmodern}

% ------------------------------------------------------------
% MATEMÁTICA E UNIDADES
% ------------------------------------------------------------

\usepackage{amsmath,amssymb,mathtools}
\usepackage{siunitx}

\sisetup{
	output-decimal-marker={,},
	per-mode=symbol,
	exponent-product=\cdot
}

% ------------------------------------------------------------
% PÁGINA E TABELAS
% ------------------------------------------------------------

\usepackage{geometry}

\geometry{
	top=1.6cm,
	bottom=1.7cm,
	left=1.7cm,
	right=1.7cm
}

\usepackage{array}
\usepackage{booktabs}
\usepackage{tabularx}
\usepackage{multirow}
\usepackage{float}
\usepackage{enumitem}

% ------------------------------------------------------------
% GRÁFICOS E TIKZ
% ------------------------------------------------------------

\usepackage{tikz}
\usepackage{pgfplots}

\pgfplotsset{compat=1.18}

\usetikzlibrary{
	arrows.meta,
	positioning,
	calc,
	patterns,
	decorations.pathreplacing,
	shapes.geometric
}

% ------------------------------------------------------------
% ASPETO GRÁFICO
% ------------------------------------------------------------

\usepackage{xcolor}
\usepackage[most]{tcolorbox}
\usepackage{fancyhdr}
\usepackage{lastpage}
\usepackage{titlesec}

% ------------------------------------------------------------
% CORES EUFISICA
% ------------------------------------------------------------

\definecolor{eufblue}{RGB}{23,86,145}
\definecolor{euflight}{RGB}{232,243,252}
\definecolor{eufdark}{RGB}{25,49,73}
\definecolor{euforange}{RGB}{230,137,47}
\definecolor{eufgreen}{RGB}{55,137,91}
\definecolor{eufred}{RGB}{190,60,55}
\definecolor{eufpurple}{RGB}{104,78,145}
\definecolor{softgray}{RGB}{246,247,249}

% ------------------------------------------------------------
% MOSTRAR / OCULTAR SOLUÇÕES
% ------------------------------------------------------------

\newif\ifsolucoes

% Para obter a versão do professor:
\solucoestrue

% Para obter apenas a ficha dos alunos, substituir por:
% \solucoesfalse

% ------------------------------------------------------------
% CABEÇALHO
% ------------------------------------------------------------

\pagestyle{fancy}
\fancyhf{}

\lhead{\textbf{\textcolor{eufblue}{EUFISICA}}}
\rhead{\textbf{Física e Química A -- 11.º ano}}

\cfoot{
	\small
	Prof. José Gonçalves
	\hfill
	Página \thepage\ de \pageref{LastPage}
}

\setlength{\headheight}{14pt}

% ------------------------------------------------------------
% FORMATAÇÃO
% ------------------------------------------------------------

\setlength{\parindent}{0pt}
\setlength{\parskip}{4pt}

\setlist[enumerate]{
	itemsep=5pt,
	topsep=4pt
}

\titleformat{\section}
{\large\bfseries\color{eufblue}}
{\thesection.}
{0.5em}
{}

% ------------------------------------------------------------
% CAIXAS
% ------------------------------------------------------------

\newtcolorbox{objetivos}{
	colback=euflight,
	colframe=eufblue,
	boxrule=0.9pt,
	arc=2mm,
	title=\textbf{Objetivos de aprendizagem},
	fonttitle=\bfseries
}

\newtcolorbox{recorda}{
	colback=softgray,
	colframe=eufdark,
	boxrule=0.7pt,
	arc=2mm,
	title=\textbf{Recorda},
	fonttitle=\bfseries
}

\newtcolorbox{dados}{
	colback=eufgreen!5,
	colframe=eufgreen,
	boxrule=0.8pt,
	arc=2mm,
	title=\textbf{Dados e constantes},
	fonttitle=\bfseries
}

\newtcolorbox{atencao}{
	colback=euforange!8,
	colframe=euforange,
	boxrule=0.8pt,
	arc=2mm,
	title=\textbf{Atenção},
	fonttitle=\bfseries
}

\newtcolorbox{stembox}{
	colback=eufpurple!5,
	colframe=eufpurple,
	boxrule=0.8pt,
	arc=2mm,
	title=\textbf{Desafio STEM},
	fonttitle=\bfseries
}

\newtcolorbox{solucaobox}[1][]{
	colback=eufgreen!4,
	colframe=eufgreen!80!black,
	boxrule=0.7pt,
	arc=2mm,
	title=\textbf{#1},
	fonttitle=\bfseries
}

\newcommand{\resposta}[1][1]{
	\vspace{#1cm}
}

% ============================================================
% DOCUMENTO
% ============================================================

\begin{document}
	
	% ============================================================
	% CAPA DA FICHA
	% ============================================================
	
	\begin{center}
		
		{\Huge\bfseries\textcolor{eufblue}{EUFISICA}}
		
		\vspace{1mm}
		
		{\Large\bfseries Ficha de Trabalho}
		
		\vspace{1mm}
		
		{\LARGE\bfseries Lei da Gravitação Universal}
		
		\vspace{2mm}
		
		{\large Física e Química A -- 11.º ano}
		
		\vspace{2mm}
		
		{\small Prof. José Gonçalves}
		
	\end{center}
	
	\vspace{2mm}
	
	\begin{tabularx}{\textwidth}{|X|X|}
		\hline
		\textbf{Nome:} & \textbf{N.º:} \\[0.45cm]
		\hline
		\textbf{Turma:} & \textbf{Data:} \\[0.45cm]
		\hline
	\end{tabularx}
	
	\vspace{4mm}
	
	% ============================================================
	% OBJETIVOS
	% ============================================================
	
	\begin{objetivos}
		
		No final desta ficha deverás ser capaz de:
		
		\begin{itemize}
			
			\item interpretar a Lei da Gravitação Universal de Newton;
			
			\item calcular a intensidade da força gravítica entre dois corpos;
			
			\item identificar a direção e o sentido das forças gravíticas;
			
			\item relacionar a força gravítica com as massas dos corpos e com a
			distância entre eles;
			
			\item relacionar a força gravítica com o peso de um corpo;
			
			\item calcular a aceleração gravítica à superfície de um planeta e a
			diferentes altitudes;
			
			\item interpretar gráficos associados à dependência
			$F_g\propto 1/r^2$;
			
			\item utilizar conhecimentos de massa volúmica, quantidade de matéria,
			massa molar e energia adquiridos no 10.º ano;
			
			\item aplicar a gravitação a situações relacionadas com planetas,
			satélites, asteroides e exploração espacial.
			
		\end{itemize}
		
	\end{objetivos}
	
	% ============================================================
	% DADOS
	% ============================================================
	
	\begin{dados}
		
		Sempre que necessário, utiliza:
		
		\[
		G=
		\SI{6.67430e-11}
		{\newton\metre\squared\per\kilogram\squared}
		\]
		
		\[
		M_T=\SI{5.972e24}{\kilogram}
		\]
		
		\[
		R_T=\SI{6.371e6}{\metre}
		\]
		
		\[
		M_L=\SI{7.342e22}{\kilogram}
		\]
		
		\[
		d_{T-L}=\SI{3.844e8}{\metre}
		\]
		
		\[
		N_A=\SI{6.022e23}{\per\mole}
		\]
		
		Considera, quando adequado,
		
		\[
		g\simeq\SI{9.8}{\metre\per\second\squared}.
		\]
		
	\end{dados}
	
	% ============================================================
	\section{Gravitação: uma interação universal}
	% ============================================================
	
	Lê o texto seguinte.
	
	\begin{tcolorbox}[
		colback=euflight,
		colframe=eufblue
		]
		
		Newton propôs que a interação que faz uma maçã cair para a superfície
		da Terra é da mesma natureza que a interação que mantém a Lua em órbita
		da Terra e os planetas em órbita do Sol.
		
		A gravitação é uma interação universal: quaisquer dois corpos com massa
		atraem-se mutuamente. A intensidade dessa interação depende das massas
		dos corpos e da distância que separa os seus centros.
		
		Mesmo quando um astronauta parece ``flutuar'' numa estação espacial,
		a atração gravítica da Terra continua a atuar. A sensação de ausência de
		peso resulta do estado de queda livre contínua da estação e dos seus
		ocupantes em torno da Terra.
		
	\end{tcolorbox}
	
	\begin{enumerate}[label=\textbf{\arabic*.}]
		
		\item Explica o significado da expressão
		\textit{``a gravitação é uma interação universal''}.
		
		\resposta[0.9]
		
		\item A força gravítica só existe entre astros?
		
		Justifica.
		
		\resposta[0.9]
		
		\item Explica por que razão não é correto afirmar que um astronauta em
		órbita está numa região onde não existe gravidade.
		
		\resposta[1.0]
		
		\item Indica dois fatores dos quais depende a intensidade da interação
		gravítica.
		
		\resposta[0.7]
		
	\end{enumerate}
	
	% ============================================================
	\section{A Lei da Gravitação Universal}
	% ============================================================
	
	\begin{recorda}
		
		A intensidade da força gravítica entre dois corpos de massas $m_1$ e
		$m_2$, cujos centros se encontram separados por uma distância $r$, é:
		
		\[
		\boxed{
			F_g=
			G\frac{m_1m_2}{r^2}
		}
		\]
		
		A força gravítica:
		
		\begin{itemize}
			
			\item tem a direção da reta que une os centros dos dois corpos;
			
			\item é sempre atrativa;
			
			\item constitui um par ação--reação;
			
			\item apresenta a mesma intensidade em ambos os corpos.
			
		\end{itemize}
		
	\end{recorda}
	
	Observa a representação seguinte.
	
	\begin{figure}[H]
		
		\centering
		
		\begin{tikzpicture}[>=Stealth,scale=1.05]
			
			% Terra
			\shade[
			ball color=eufblue!70
			]
			(-4,0) circle (1.25);
			
			\node[white,font=\bfseries] at (-4,0) {Terra};
			
			% Lua
			\shade[
			ball color=gray!60
			]
			(4,0) circle (0.65);
			
			\node[font=\small] at (4,0) {Lua};
			
			% centros
			\fill (-4,0) circle (0.05);
			\fill (4,0) circle (0.05);
			
			% distância
			\draw[
			<->,
			thick
			]
			(-4,-1.7) -- (4,-1.7);
			
			\node[
			fill=white,
			inner sep=2pt
			]
			at (0,-1.7)
			{$r$};
			
			% forças
			\draw[
			->,
			very thick,
			eufred
			]
			(-2.7,0.4) -- (-0.5,0.4);
			
			\node[
			above,
			text=eufred
			]
			at (-1.6,0.4)
			{$\vec F_{L\rightarrow T}$};
			
			\draw[
			->,
			very thick,
			eufgreen!70!black
			]
			(3.3,0.4) -- (1.1,0.4);
			
			\node[
			above,
			text=eufgreen!70!black
			]
			at (2.2,0.4)
			{$\vec F_{T\rightarrow L}$};
			
		\end{tikzpicture}
		
		\caption{Interação gravítica entre a Terra e a Lua.}
		
	\end{figure}
	
	\begin{enumerate}[label=\textbf{\arabic*.},resume]
		
		\item Indica a direção da força gravítica exercida pela Terra sobre a Lua.
		
		\resposta[0.6]
		
		\item Indica o sentido da força gravítica exercida pela Terra sobre a Lua.
		
		\resposta[0.6]
		
		\item A Terra tem uma massa muito superior à da Lua.
		
		Será a força exercida pela Terra sobre a Lua maior do que a força
		exercida pela Lua sobre a Terra?
		
		Justifica.
		
		\resposta[1.0]
		
		\item Representa, através de vetores, as duas forças gravíticas na
		figura seguinte.
		
		\begin{center}
			
			\begin{tikzpicture}[>=Stealth]
				
				\shade[ball color=euforange!75]
				(-3,0) circle (0.8);
				
				\shade[ball color=eufblue!65]
				(3,0) circle (1.1);
				
				\node at (-3,-1.2) {$m_1$};
				\node at (3,-1.5) {$m_2$};
				
				\draw[dashed]
				(-3,0)--(3,0);
				
			\end{tikzpicture}
			
		\end{center}
		
	\end{enumerate}
	
	% ============================================================
	\section{Do peso à Lei da Gravitação Universal}
	% ============================================================
	
	Uma pessoa de massa \SI{65}{\kilogram} encontra-se à superfície da Terra.
	
	Considera que a distância entre a pessoa e o centro da Terra é
	aproximadamente igual ao raio terrestre.
	
	\begin{enumerate}[label=\textbf{\arabic*.},resume]
		
		\item Utilizando a Lei da Gravitação Universal, calcula a intensidade da
		força gravítica exercida pela Terra sobre a pessoa.
		
		\resposta[1.5]
		
		\item Calcula o valor da aceleração gravítica à superfície da Terra a partir
		da expressão
		
		\[
		g=\frac{F_g}{m}.
		\]
		
		\resposta[1.2]
		
		\item Calcula o peso da pessoa utilizando
		
		\[
		P=mg,
		\]
		
		com $g=\SI{9.8}{\metre\per\second\squared}$.
		
		\resposta[0.9]
		
		\item Compara os resultados das questões 9 e 11.
		
		O que podes concluir?
		
		\resposta[1.0]
		
		\item Mostra que, junto à superfície da Terra,
		
		\[
		\boxed{
			g=
			G\frac{M_T}{R_T^2}
		}
		\]
		
		\resposta[1.2]
		
	\end{enumerate}
	
	% ============================================================
	\section{A interação Terra--Lua}
	% ============================================================
	
	Considera os dados fornecidos no início da ficha.
	
	\begin{enumerate}[label=\textbf{\arabic*.},resume]
		
		\item Calcula a intensidade da força gravítica exercida pela Terra sobre
		a Lua.
		
		\resposta[1.5]
		
		\item Qual é a intensidade da força exercida pela Lua sobre a Terra?
		
		Justifica.
		
		\resposta[0.8]
		
		\item Imagina que a distância entre a Terra e a Lua duplicava, mantendo-se
		as massas constantes.
		
		Qual seria a nova intensidade da força gravítica?
		
		\resposta[1.0]
		
		\item Se a massa da Lua fosse metade da atual, mantendo-se todas as outras
		grandezas constantes, qual seria a intensidade da força gravítica?
		
		\resposta[1.0]
		
		\item Completa:
		
		\[
		F_g\propto \underline{\hspace{2cm}}
		\]
		
		relativamente às massas e
		
		\[
		F_g\propto \underline{\hspace{2cm}}
		\]
		
		relativamente à distância.
		
	\end{enumerate}
	
	% ============================================================
	\section{Como varia a força com a distância?}
	% ============================================================
	
	Para dois corpos de massas constantes, considera $F_0$ a força gravítica
	quando a distância entre os seus centros é $r_0$.
	
	O gráfico seguinte representa a razão $F_g/F_0$ em função de $r/r_0$.
	
	\begin{figure}[H]
		
		\centering
		
		\begin{tikzpicture}
			
			\begin{axis}[
				width=12cm,
				height=7cm,
				axis lines=left,
				xlabel={$r/r_0$},
				ylabel={$F_g/F_0$},
				xmin=0.8,
				xmax=4.1,
				ymin=0,
				ymax=1.1,
				xtick={1,1.5,2,2.5,3,3.5,4},
				ytick={0,0.1,0.2,0.3,0.4,0.5,0.6,0.7,0.8,0.9,1},
				grid=both,
				major grid style={gray!30}
				]
				
				\addplot[
				domain=1:4,
				samples=150,
				very thick,
				eufblue
				]
				{1/(x^2)};
				
			\end{axis}
			
		\end{tikzpicture}
		
		\caption{Dependência da força gravítica com a distância.}
		
	\end{figure}
	
	\begin{enumerate}[label=\textbf{\arabic*.},resume]
		
		\item Quando a distância duplica, que fração do valor inicial assume
		a força gravítica?
		
		\resposta[0.6]
		
		\item E quando a distância triplica?
		
		\resposta[0.6]
		
		\item Se a distância passar de $r_0$ para $4r_0$, qual será a razão
		$F_g/F_0$?
		
		\resposta[0.7]
		
		\item Completa a tabela.
		
		\begin{table}[H]
			
			\centering
			
			\renewcommand{\arraystretch}{1.4}
			
			\begin{tabular}{c|ccccc}
				
				\toprule
				
				$r/r_0$
				&
				1,0
				&
				1,5
				&
				2,0
				&
				2,5
				&
				3,0
				
				\\
				
				\midrule
				
				$F_g/F_0$
				&
				1,000
				&
				&
				&
				&
				\\
				
				\bottomrule
				
			\end{tabular}
			
		\end{table}
		
		\item Utiliza os valores obtidos para representar os pontos no gráfico
		seguinte.
		
		\begin{center}
			
			\begin{tikzpicture}
				
				\begin{axis}[
					width=12cm,
					height=7cm,
					axis lines=left,
					xlabel={$r/r_0$},
					ylabel={$F_g/F_0$},
					xmin=0,
					xmax=3.5,
					ymin=0,
					ymax=1.1,
					xtick={0,0.5,1,1.5,2,2.5,3,3.5},
					ytick={0,0.1,0.2,0.3,0.4,0.5,0.6,0.7,0.8,0.9,1},
					grid=both,
					major grid style={gray!30}
					]
					
				\end{axis}
				
			\end{tikzpicture}
			
		\end{center}
		
		\item Explica por que razão o gráfico não é uma reta.
		
		\resposta[0.9]
		
	\end{enumerate}
	
	% ============================================================
	\section{Um satélite em órbita terrestre}
	% ============================================================
	
	Um satélite de massa
	
	\[
	m=\SI{1200}{\kilogram}
	\]
	
	encontra-se numa órbita circular a uma altitude de
	
	\[
	h=\SI{400}{\kilo\metre}
	\]
	
	acima da superfície terrestre.
	
	\begin{center}
		
		\begin{tikzpicture}[scale=0.95]
			
			\shade[
			ball color=eufblue!70
			]
			(0,0) circle (1.5);
			
			\node[
			white,
			font=\bfseries
			]
			at (0,0)
			{Terra};
			
			\draw[
			dashed,
			eufdark
			]
			(0,0) circle (3.2);
			
			\filldraw[
			fill=euforange,
			draw=euforange!70!black
			]
			(3.2,0) circle (0.18);
			
			\node[
			right
			]
			at (3.4,0)
			{satélite};
			
			\draw[
			->,
			very thick,
			eufred
			]
			(3.1,0.2)--(1.7,0.2);
			
			\node[
			above
			]
			at (2.35,0.2)
			{$\vec F_g$};
			
			\draw[
			<->,
			thick
			]
			(0,-0.4)--(3.2,-0.4);
			
			\node[
			fill=white,
			inner sep=2pt
			]
			at (1.6,-0.4)
			{$r=R_T+h$};
			
		\end{tikzpicture}
		
	\end{center}
	
	\begin{enumerate}[label=\textbf{\arabic*.},resume]
		
		\item Calcula a distância do satélite ao centro da Terra.
		
		\resposta[0.8]
		
		\item Calcula a intensidade da força gravítica exercida pela Terra sobre
		o satélite.
		
		\resposta[1.3]
		
		\item Calcula a aceleração gravítica à altitude do satélite.
		
		\resposta[1.2]
		
		\item Compara esse valor com o valor da aceleração gravítica à superfície
		da Terra.
		
		\resposta[1.0]
		
		\item Um aluno afirma:
		
		\begin{quote}
			``A \SI{400}{\kilo\metre} de altitude praticamente já não existe
			gravidade.''
		\end{quote}
		
		Comenta a afirmação utilizando os cálculos efetuados.
		
		\resposta[1.1]
		
	\end{enumerate}
	
	% ============================================================
	\section{Determinar a massa de um planeta}
	% ============================================================
	
	À superfície de Marte, a aceleração gravítica é aproximadamente
	
	\[
	g_M=\SI{3.71}{\metre\per\second\squared}.
	\]
	
	O raio médio de Marte é
	
	\[
	R_M=\SI{3.39e6}{\metre}.
	\]
	
	\begin{enumerate}[label=\textbf{\arabic*.},resume]
		
		\item Partindo de
		
		\[
		g=
		G\frac{M}{R^2},
		\]
		
		deduz uma expressão para determinar a massa $M$ de um planeta.
		
		\resposta[0.8]
		
		\item Calcula a massa de Marte.
		
		\resposta[1.3]
		
		\item Determina o peso que um astronauta de massa \SI{70}{\kilogram}
		teria à superfície de Marte.
		
		\resposta[0.9]
		
		\item A massa do astronauta seria diferente da sua massa na Terra?
		
		Explica.
		
		\resposta[0.8]
		
	\end{enumerate}
	
	% ============================================================
	\section{Aceleração gravítica e altitude}
	% ============================================================
	
	A aceleração gravítica a uma distância $r$ do centro da Terra é
	
	\[
	g(r)=G\frac{M_T}{r^2}.
	\]
	
	\begin{enumerate}[label=\textbf{\arabic*.},resume]
		
		\item A que distância do centro da Terra a aceleração gravítica será
		igual a um quarto do seu valor à superfície?
		
		\resposta[1.3]
		
		\item Qual será, nesse caso, a altitude relativamente à superfície da Terra?
		
		\resposta[1.0]
		
		\item Se a distância ao centro da Terra triplicar relativamente ao raio
		terrestre, qual será aproximadamente o valor de $g$?
		
		\resposta[1.0]
		
	\end{enumerate}
	
	% ============================================================
	\section{Ligação ao 10.º ano -- massa volúmica de um asteroide}
	% ============================================================
	
	Um pequeno asteroide pode ser aproximadamente modelado como uma esfera
	de raio
	
	\[
	R_A=\SI{250}{\metre}.
	\]
	
	Admite que é constituído essencialmente por material rochoso e metálico,
	com massa volúmica média
	
	\[
	\rho=\SI{7800}{\kilogram\per\metre\cubed}.
	\]
	
	Recorda que:
	
	\[
	\rho=\frac{m}{V}
	\]
	
	e que o volume de uma esfera é
	
	\[
	V=\frac{4}{3}\pi R^3.
	\]
	
	Uma sonda de massa \SI{1000}{\kilogram} encontra-se a uma distância de
	\SI{1000}{\metre} do centro do asteroide.
	
	\begin{center}
		
		\begin{tikzpicture}[>=Stealth]
			
			\shade[
			ball color=gray!65
			]
			(-3,0) circle (1.1);
			
			\node at (-3,-1.5)
			{Asteroide};
			
			\filldraw[
			fill=euforange,
			draw=euforange!70!black
			]
			(3,0) rectangle +(0.45,0.28);
			
			\node at (3.25,-0.55)
			{sonda};
			
			\draw[
			<->,
			thick
			]
			(-3,-2)--(3.25,-2);
			
			\node[
			fill=white
			]
			at (0.1,-2)
			{\SI{1000}{m}};
			
			\draw[
			->,
			very thick,
			eufred
			]
			(3,0.2)--(1.5,0.2);
			
			\node[
			above
			]
			at (2.2,0.2)
			{$\vec F_g$};
			
		\end{tikzpicture}
		
	\end{center}
	
	\begin{enumerate}[label=\textbf{\arabic*.},resume]
		
		\item Calcula o volume aproximado do asteroide.
		
		\resposta[1.0]
		
		\item Utilizando a massa volúmica, calcula a massa aproximada do asteroide.
		
		\resposta[1.2]
		
		\item Calcula a intensidade da força gravítica exercida pelo asteroide
		sobre a sonda.
		
		\resposta[1.3]
		
		\item Compara a força calculada com o peso da mesma sonda à superfície
		da Terra.
		
		\resposta[1.0]
		
		\item Explica por que razão um corpo relativamente grande pode apresentar
		uma atração gravítica muito inferior à da Terra.
		
		\resposta[1.0]
		
	\end{enumerate}
	
	% ============================================================
	\section{Ligação à Química do 10.º ano -- quantidade de matéria}
	% ============================================================
	
	Considera duas amostras de carbono puro, cada uma contendo
	
	\[
	n=\SI{2.00}{\mole}
	\]
	
	de átomos de carbono.
	
	Admite, num modelo muito simplificado, que as duas amostras podem ser
	tratadas como massas puntiformes separadas por
	
	\[
	r=\SI{0.10}{\metre}.
	\]
	
	Utiliza
	
	\[
	M(C)=\SI{12.0}{\gram\per\mole}.
	\]
	
	\begin{enumerate}[label=\textbf{\arabic*.},resume]
		
		\item Calcula a massa de cada amostra através da relação
		
		\[
		m=nM.
		\]
		
		Apresenta o resultado em quilogramas.
		
		\resposta[1.0]
		
		\item Calcula o número de átomos existente em cada amostra.
		
		\resposta[0.9]
		
		\item Calcula a intensidade da força gravítica entre as duas amostras.
		
		\resposta[1.3]
		
		\item Explica por que razão a interação gravítica entre objetos de pequena
		massa é geralmente impercetível em laboratório.
		
		\resposta[1.0]
		
		\item Se cada amostra tivesse o dobro da quantidade de matéria, mantendo-se
		a distância, por que fator seria multiplicada a força gravítica?
		
		Justifica.
		
		\resposta[0.9]
		
	\end{enumerate}
	
	% ============================================================
	\section{Uma ``Super-Terra''}
	% ============================================================
	
	Foi descoberto um planeta hipotético cuja massa é cinco vezes a massa
	da Terra e cujo raio é $1,6$ vezes o raio terrestre:
	
	\[
	M_P=5M_T
	\]
	
	e
	
	\[
	R_P=1{,}6R_T.
	\]
	
	\begin{enumerate}[label=\textbf{\arabic*.},resume]
		
		\item Sem substituir inicialmente os valores de $G$, $M_T$ e $R_T$,
		mostra que:
		
		\[
		\frac{g_P}{g_T}
		=
		\frac{5}{(1{,}6)^2}.
		\]
		
		\resposta[1.1]
		
		\item Calcula a aceleração gravítica aproximada à superfície desse planeta.
		
		\resposta[1.0]
		
		\item Determina o peso de uma pessoa de massa \SI{70}{\kilogram}
		nesse planeta.
		
		\resposta[0.9]
		
		\item A massa da pessoa altera-se quando viaja da Terra para o planeta?
		
		\resposta[0.7]
		
		\item O facto de o planeta ter cinco vezes a massa da Terra significa que
		a aceleração gravítica à sua superfície também é cinco vezes superior?
		
		Justifica.
		
		\resposta[1.0]
		
	\end{enumerate}
	
	% ============================================================
	\section{Tratamento de dados e representação gráfica}
	% ============================================================
	
	Num modelo laboratorial idealizado, consideram-se duas massas de
	
	\[
	m_1=m_2=\SI{1000}{\kilogram}.
	\]
	
	A força gravítica é calculada para várias distâncias.
	
	Completa a tabela seguinte.
	
	\begin{table}[H]
		
		\centering
		
		\renewcommand{\arraystretch}{1.45}
		
		\begin{tabular}{c|c|c}
			
			\toprule
			
			$r$/m
			&
			$1/r^2$/m$^{-2}$
			&
			$F_g$/N
			
			\\
			
			\midrule
			
			1,0 & & \\
			
			1,5 & & \\
			
			2,0 & & \\
			
			2,5 & & \\
			
			3,0 & & \\
			
			\bottomrule
			
		\end{tabular}
		
	\end{table}
	
	\begin{enumerate}[label=\textbf{\arabic*.},resume]
		
		\item Calcula os valores em falta na tabela.
		
		\resposta[1.8]
		
		\item Representa aproximadamente $F_g$ em função de $r$.
		
		\begin{center}
			
			\begin{tikzpicture}
				
				\begin{axis}[
					width=12cm,
					height=6.5cm,
					axis lines=left,
					xlabel={$r$/m},
					ylabel={$F_g$/N},
					xmin=0,
					xmax=3.5,
					ymin=0,
					ymax=0.000075,
					xtick={0,0.5,1,1.5,2,2.5,3,3.5},
					scaled y ticks=base 10:5,
					grid=both,
					major grid style={gray!30}
					]
					
				\end{axis}
				
			\end{tikzpicture}
			
		\end{center}
		
		\item O gráfico $F_g=f(r)$ deverá ser uma reta?
		
		Justifica.
		
		\resposta[0.8]
		
		\item Se representasses $F_g$ em função de $1/r^2$, que tipo de gráfico
		esperarias obter?
		
		\resposta[0.8]
		
		\item Determina o declive teórico desse gráfico $F_g=f(1/r^2)$.
		
		\resposta[1.0]
		
	\end{enumerate}
	
	% ============================================================
	\section{Desafio STEM -- um satélite em órbita circular}
	% ============================================================
	
	\begin{stembox}
		
		Um satélite de massa \SI{1200}{\kilogram} descreve uma órbita circular
		a \SI{400}{\kilo\metre} de altitude.
		
		Num movimento circular uniforme, a resultante das forças tem intensidade
		
		\[
		F_R=m\frac{v^2}{r}.
		\]
		
		No caso do satélite, admite que a força gravítica é a força resultante.
		
	\end{stembox}
	
	\begin{enumerate}[label=\textbf{\arabic*.},resume]
		
		\item Igualando
		
		\[
		G\frac{M_Tm}{r^2}
		=
		m\frac{v^2}{r},
		\]
		
		mostra que a velocidade orbital é dada por
		
		\[
		\boxed{
			v=
			\sqrt{\frac{GM_T}{r}}
		}.
		\]
		
		\resposta[1.1]
		
		\item Calcula a velocidade orbital do satélite.
		
		\resposta[1.2]
		
		\item Utilizando
		
		\[
		T=\frac{2\pi r}{v},
		\]
		
		determina o período orbital.
		
		Apresenta o resultado em minutos.
		
		\resposta[1.2]
		
		\item Recordando do 10.º ano que
		
		\[
		E_c=\frac{1}{2}mv^2,
		\]
		
		calcula a energia cinética do satélite.
		
		\resposta[1.2]
		
		\item Explica por que razão a massa do satélite não aparece na expressão
		final da velocidade orbital.
		
		\resposta[1.0]
		
	\end{enumerate}
	
	% ============================================================
	\section{Escolha múltipla}
	% ============================================================
	
	\begin{enumerate}[label=\textbf{\arabic*.},resume]
		
		\item Se uma das massas duplicar e a distância se mantiver, a força
		gravítica:
		
		\begin{enumerate}[label=\Alph*.]
			
			\item reduz-se a metade;
			
			\item duplica;
			
			\item quadruplica;
			
			\item não se altera.
			
		\end{enumerate}
		
		\item Se as duas massas duplicarem e a distância se mantiver, a força
		gravítica:
		
		\begin{enumerate}[label=\Alph*.]
			
			\item duplica;
			
			\item quadruplica;
			
			\item reduz-se a metade;
			
			\item aumenta oito vezes.
			
		\end{enumerate}
		
		\item Se a distância entre dois corpos triplicar, a força gravítica passa
		a:
		
		\begin{enumerate}[label=\Alph*.]
			
			\item $3F$;
			
			\item $F/3$;
			
			\item $F/6$;
			
			\item $F/9$.
			
		\end{enumerate}
		
		\item As forças gravíticas exercidas mutuamente por dois corpos apresentam:
		
		\begin{enumerate}[label=\Alph*.]
			
			\item intensidades diferentes porque as massas são diferentes;
			
			\item a mesma direção, o mesmo sentido e diferentes intensidades;
			
			\item a mesma intensidade e sentidos opostos;
			
			\item intensidades e direções diferentes.
			
		\end{enumerate}
		
		\item Um astronauta vai da Terra para a Lua.
		
		Qual das seguintes grandezas permanece constante?
		
		\begin{enumerate}[label=\Alph*.]
			
			\item peso;
			
			\item força gravítica;
			
			\item massa;
			
			\item aceleração gravítica.
			
		\end{enumerate}
		
		\item A aceleração gravítica criada por um planeta à distância $r$ do seu
		centro:
		
		\begin{enumerate}[label=\Alph*.]
			
			\item depende da massa do corpo colocado nesse ponto;
			
			\item aumenta proporcionalmente a $r^2$;
			
			\item é dada por $GM/r^2$;
			
			\item é sempre igual a \SI{9.8}{\metre\per\second\squared}.
			
		\end{enumerate}
		
	\end{enumerate}
	
	% ============================================================
	\section{Problema final de síntese}
	% ============================================================
	
	Um planeta esférico apresenta:
	
	\[
	M=\SI{8.0e23}{\kilogram}
	\]
	
	e
	
	\[
	R=\SI{4.0e6}{\metre}.
	\]
	
	Uma sonda de massa \SI{500}{\kilogram} encontra-se inicialmente à
	superfície do planeta.
	
	Posteriormente, sobe até uma altitude igual ao raio do planeta.
	
	\begin{enumerate}[label=\textbf{\arabic*.},resume]
		
		\item Calcula a aceleração gravítica à superfície do planeta.
		
		\resposta[1.2]
		
		\item Calcula o peso da sonda à superfície.
		
		\resposta[0.9]
		
		\item Quando a altitude é igual ao raio do planeta, qual é a distância da
		sonda ao centro do planeta?
		
		\resposta[0.7]
		
		\item Calcula a aceleração gravítica nessa altitude.
		
		\resposta[1.1]
		
		\item Calcula a força gravítica exercida sobre a sonda nessa posição.
		
		\resposta[1.0]
		
		\item Determina a razão
		
		\[
		\frac{F_{\text{altitude}}}
		{F_{\text{superfície}}}.
		\]
		
		\resposta[0.9]
		
		\item Explica o resultado da questão anterior utilizando a dependência da
		força gravítica com a distância.
		
		\resposta[1.1]
		
	\end{enumerate}
	
	% ============================================================
	% SOLUÇÕES
	% ============================================================
	
	\ifsolucoes
	
	\clearpage
	
	\begin{center}
		
		{\Huge\bfseries\textcolor{eufgreen!70!black}{Resolução}}
		
		\vspace{2mm}
		
		{\large Lei da Gravitação Universal}
		
		\vspace{1mm}
		
		{\small EUFISICA -- Prof. José Gonçalves}
		
	\end{center}
	
	\vspace{4mm}
	
	% ============================================================
	\section*{Resolução 1 -- Gravitação universal}
	% ============================================================
	
	\begin{solucaobox}[Questões 1--4]
		
		\textbf{1.}
		
		A gravitação é considerada universal porque quaisquer dois corpos que
		tenham massa exercem forças gravíticas um sobre o outro.
		
		\vspace{2mm}
		
		\textbf{2.}
		
		Não. Todos os corpos com massa interagem graviticamente.
		
		A força é muito mais evidente entre astros porque estes apresentam massas
		muito elevadas.
		
		\vspace{2mm}
		
		\textbf{3.}
		
		Um astronauta em órbita continua sujeito à força gravítica terrestre.
		
		A estação espacial e o astronauta encontram-se em queda livre em torno da
		Terra, o que origina a sensação aparente de ausência de peso.
		
		\vspace{2mm}
		
		\textbf{4.}
		
		A intensidade da força gravítica depende:
		
		\begin{itemize}
			
			\item das massas dos dois corpos;
			
			\item da distância entre os respetivos centros.
			
		\end{itemize}
		
	\end{solucaobox}
	
	% ============================================================
	\section*{Resolução 2 -- Direção, sentido e par ação--reação}
	% ============================================================
	
	\begin{solucaobox}[Questões 5--8]
		
		\textbf{5.}
		
		A direção é a da reta que une os centros da Terra e da Lua.
		
		\vspace{2mm}
		
		\textbf{6.}
		
		A força exercida pela Terra sobre a Lua aponta da Lua para o centro da
		Terra.
		
		\vspace{2mm}
		
		\textbf{7.}
		
		Não.
		
		De acordo com a terceira lei de Newton,
		
		\[
		\left|
		\vec F_{T\rightarrow L}
		\right|
		=
		\left|
		\vec F_{L\rightarrow T}
		\right|.
		\]
		
		As forças apresentam a mesma intensidade e sentidos opostos.
		
		\vspace{2mm}
		
		\textbf{8.}
		
		Os dois vetores devem:
		
		\begin{itemize}
			
			\item estar sobre a linha que une os centros;
			
			\item apontar um para o outro;
			
			\item apresentar o mesmo comprimento, representando igual intensidade.
			
		\end{itemize}
		
	\end{solucaobox}
	
	% ============================================================
	\section*{Resolução 3 -- Peso e gravitação}
	% ============================================================
	
	\begin{solucaobox}[Questões 9--13]
		
		\textbf{9.}
		
		\[
		F_g=
		G\frac{M_Tm}{R_T^2}
		\]
		
		\[
		F_g=
		\frac{
			(6{,}67430\times10^{-11})
			(5{,}972\times10^{24})
			(65)
		}{
			(6{,}371\times10^6)^2
		}
		\]
		
		\[
		\boxed{
			F_g\approx\SI{6.38e2}{N}
		}
		\]
		
		ou aproximadamente
		
		\[
		\boxed{
			F_g\approx\SI{638}{N}
		}
		\]
		
		\vspace{3mm}
		
		\textbf{10.}
		
		\[
		g=\frac{F_g}{m}
		\]
		
		\[
		g=
		\frac{638}{65}
		\]
		
		\[
		\boxed{
			g\approx\SI{9.82}{m.s^{-2}}
		}
		\]
		
		\vspace{3mm}
		
		\textbf{11.}
		
		\[
		P=mg
		\]
		
		\[
		P=65\times9{,}8
		\]
		
		\[
		\boxed{
			P\approx\SI{637}{N}
		}
		\]
		
		\vspace{3mm}
		
		\textbf{12.}
		
		Os valores são praticamente iguais.
		
		A pequena diferença resulta das aproximações utilizadas.
		
		Conclui-se que o peso de um corpo junto à superfície terrestre corresponde
		à força gravítica exercida pela Terra sobre o corpo.
		
		\vspace{3mm}
		
		\textbf{13.}
		
		Da Lei da Gravitação Universal,
		
		\[
		F_g=
		G\frac{M_Tm}{R_T^2}.
		\]
		
		Como
		
		\[
		F_g=mg,
		\]
		
		então
		
		\[
		mg=
		G\frac{M_Tm}{R_T^2}.
		\]
		
		Cancelando $m$:
		
		\[
		\boxed{
			g=
			G\frac{M_T}{R_T^2}
		}.
		\]
		
		Isto mostra também que, desprezando a resistência do ar, a aceleração
		gravítica não depende da massa do corpo em queda.
		
	\end{solucaobox}
	
	% ============================================================
	\section*{Resolução 4 -- Terra e Lua}
	% ============================================================
	
	\begin{solucaobox}[Questões 14--18]
		
		\textbf{14.}
		
		\[
		F_g=
		G\frac{M_TM_L}{r^2}
		\]
		
		\[
		F_g=
		\frac{
			(6{,}67430\times10^{-11})
			(5{,}972\times10^{24})
			(7{,}342\times10^{22})
		}{
			(3{,}844\times10^8)^2
		}
		\]
		
		\[
		\boxed{
			F_g\approx\SI{1.98e20}{N}
		}
		\]
		
		\vspace{3mm}
		
		\textbf{15.}
		
		A intensidade da força da Lua sobre a Terra é também
		
		\[
		\boxed{
			\SI{1.98e20}{N}
		}
		\]
		
		pela terceira lei de Newton.
		
		\vspace{3mm}
		
		\textbf{16.}
		
		Se a distância duplicar:
		
		\[
		F'_g=
		\frac{F_g}{2^2}
		=
		\frac{F_g}{4}.
		\]
		
		Logo,
		
		\[
		F'_g=
		\frac{1{,}98\times10^{20}}{4}
		\]
		
		\[
		\boxed{
			F'_g\approx\SI{4.95e19}{N}
		}
		\]
		
		\vspace{3mm}
		
		\textbf{17.}
		
		Como
		
		\[
		F_g\propto M_L,
		\]
		
		reduzir a massa da Lua para metade reduz a força para metade:
		
		\[
		F'_g=
		\frac{1{,}98\times10^{20}}{2}
		\]
		
		\[
		\boxed{
			F'_g\approx\SI{9.90e19}{N}
		}
		\]
		
		\vspace{3mm}
		
		\textbf{18.}
		
		\[
		\boxed{
			F_g\propto m_1m_2
		}
		\]
		
		e
		
		\[
		\boxed{
			F_g\propto\frac{1}{r^2}.
		}
		\]
		
	\end{solucaobox}
	
	% ============================================================
	\section*{Resolução 5 -- Dependência com a distância}
	% ============================================================
	
	\begin{solucaobox}[Questões 19--24]
		
		Como
		
		\[
		\frac{F_g}{F_0}
		=
		\frac{r_0^2}{r^2},
		\]
		
		tem-se:
		
		\vspace{2mm}
		
		\textbf{19.}
		
		Para
		
		\[
		r=2r_0,
		\]
		
		\[
		\frac{F_g}{F_0}
		=
		\frac{1}{4}.
		\]
		
		\vspace{2mm}
		
		\textbf{20.}
		
		Para
		
		\[
		r=3r_0,
		\]
		
		\[
		\boxed{
			\frac{F_g}{F_0}
			=
			\frac{1}{9}
		}
		\]
		
		\vspace{2mm}
		
		\textbf{21.}
		
		Para
		
		\[
		r=4r_0,
		\]
		
		\[
		\boxed{
			\frac{F_g}{F_0}
			=
			\frac{1}{16}
			=
			0{,}0625
		}
		\]
		
		\vspace{3mm}
		
		\textbf{22.}
		
		\[
		\begin{array}{c|ccccc}
			r/r_0 & 1{,}0 & 1{,}5 & 2{,}0 & 2{,}5 & 3{,}0\\
			\hline
			F_g/F_0 &
			1{,}000 &
			0{,}444 &
			0{,}250 &
			0{,}160 &
			0{,}111
		\end{array}
		\]
		
		\vspace{3mm}
		
		\textbf{23.}
		
		Os pontos devem localizar-se aproximadamente nas coordenadas
		
		\[
		(1;1),
		\quad
		(1{,}5;0{,}444),
		\quad
		(2;0{,}25),
		\]
		
		\[
		(2{,}5;0{,}16),
		\quad
		(3;0{,}111).
		\]
		
		\vspace{3mm}
		
		\textbf{24.}
		
		O gráfico não é uma reta porque a força não é inversamente proporcional
		a $r$, mas ao quadrado da distância:
		
		\[
		F_g\propto\frac{1}{r^2}.
		\]
		
	\end{solucaobox}
	
	% ============================================================
	\section*{Resolução 6 -- Satélite a 400 km}
	% ============================================================
	
	\begin{solucaobox}[Questões 25--29]
		
		\textbf{25.}
		
		Primeiro:
		
		\[
		h=\SI{400}{km}
		=
		\SI{4.00e5}{m}.
		\]
		
		Logo,
		
		\[
		r=R_T+h
		\]
		
		\[
		r=
		6{,}371\times10^6
		+
		4{,}00\times10^5
		\]
		
		\[
		\boxed{
			r=
			\SI{6.771e6}{m}
		}
		\]
		
		\vspace{3mm}
		
		\textbf{26.}
		
		\[
		F_g=
		G\frac{M_Tm}{r^2}
		\]
		
		\[
		F_g=
		\frac{
			(6{,}67430\times10^{-11})
			(5{,}972\times10^{24})
			(1200)
		}{
			(6{,}771\times10^6)^2
		}
		\]
		
		\[
		\boxed{
			F_g\approx\SI{1.04e4}{N}
		}
		\]
		
		\vspace{3mm}
		
		\textbf{27.}
		
		\[
		g=
		\frac{F_g}{m}
		\]
		
		\[
		g=
		\frac{1{,}043\times10^4}{1200}
		\]
		
		\[
		\boxed{
			g\approx\SI{8.69}{m.s^{-2}}
		}
		\]
		
		\vspace{3mm}
		
		\textbf{28.}
		
		À superfície:
		
		\[
		g_T\approx\SI{9.82}{m.s^{-2}}.
		\]
		
		Assim,
		
		\[
		\frac{g_{400}}{g_T}
		=
		\frac{8{,}69}{9{,}82}
		\approx0{,}885.
		\]
		
		Portanto, a \SI{400}{km} de altitude, $g$ ainda corresponde a cerca de
		
		\[
		\boxed{88{,}5\%}
		\]
		
		do valor à superfície.
		
		\vspace{3mm}
		
		\textbf{29.}
		
		A afirmação é falsa.
		
		A aceleração gravítica ainda é aproximadamente
		
		\[
		\SI{8.69}{m.s^{-2}},
		\]
		
		um valor bastante próximo do valor à superfície.
		
		A sensação de ausência de peso numa estação espacial deve-se à queda
		livre orbital e não à ausência de gravidade.
		
	\end{solucaobox}
	
	% ============================================================
	\section*{Resolução 7 -- Massa de Marte}
	% ============================================================
	
	\begin{solucaobox}[Questões 30--33]
		
		\textbf{30.}
		
		\[
		g=
		G\frac{M}{R^2}
		\]
		
		Logo,
		
		\[
		\boxed{
			M=
			\frac{gR^2}{G}
		}
		\]
		
		\vspace{3mm}
		
		\textbf{31.}
		
		\[
		M=
		\frac{
			3{,}71(3{,}39\times10^6)^2
		}{
			6{,}67430\times10^{-11}
		}
		\]
		
		\[
		\boxed{
			M_M\approx\SI{6.39e23}{kg}
		}
		\]
		
		\vspace{3mm}
		
		\textbf{32.}
		
		\[
		P=mg
		\]
		
		\[
		P=70\times3{,}71
		\]
		
		\[
		\boxed{
			P\approx\SI{260}{N}
		}
		\]
		
		\vspace{3mm}
		
		\textbf{33.}
		
		Não.
		
		A massa é uma propriedade do corpo e mantém-se aproximadamente constante.
		
		O peso altera-se porque depende da aceleração gravítica local.
		
	\end{solucaobox}
	
	% ============================================================
	\section*{Resolução 8 -- Gravidade e altitude}
	% ============================================================
	
	\begin{solucaobox}[Questões 34--36]
		
		\textbf{34.}
		
		À superfície:
		
		\[
		g_T=
		G\frac{M_T}{R_T^2}.
		\]
		
		Pretende-se:
		
		\[
		g'=\frac{g_T}{4}.
		\]
		
		Assim,
		
		\[
		G\frac{M_T}{r^2}
		=
		\frac14
		G\frac{M_T}{R_T^2}.
		\]
		
		Cancelando os fatores comuns:
		
		\[
		\frac1{r^2}
		=
		\frac1{4R_T^2}.
		\]
		
		Logo,
		
		\[
		\boxed{
			r=2R_T
		}
		\]
		
		ou
		
		\[
		r=
		2(6{,}371\times10^6)
		\]
		
		\[
		\boxed{
			r=
			\SI{1.2742e7}{m}
		}
		\]
		
		\vspace{3mm}
		
		\textbf{35.}
		
		Como
		
		\[
		h=r-R_T,
		\]
		
		\[
		h=2R_T-R_T=R_T.
		\]
		
		Portanto,
		
		\[
		\boxed{
			h\approx\SI{6.371e6}{m}
			=
			\SI{6371}{km}
		}
		\]
		
		\vspace{3mm}
		
		\textbf{36.}
		
		Se
		
		\[
		r=3R_T,
		\]
		
		então
		
		\[
		g=
		\frac{g_T}{3^2}
		=
		\frac{9{,}8}{9}.
		\]
		
		\[
		\boxed{
			g\approx\SI{1.09}{m.s^{-2}}
		}
		\]
		
	\end{solucaobox}
	
	% ============================================================
	\section*{Resolução 9 -- Asteroide e massa volúmica}
	% ============================================================
	
	\begin{solucaobox}[Questões 37--41]
		
		\textbf{37.}
		
		\[
		V=
		\frac43\pi R^3
		\]
		
		\[
		V=
		\frac43\pi(250)^3
		\]
		
		\[
		\boxed{
			V\approx\SI{6.54e7}{m^3}
		}
		\]
		
		\vspace{3mm}
		
		\textbf{38.}
		
		Da definição de massa volúmica:
		
		\[
		\rho=\frac{m}{V}
		\]
		
		logo,
		
		\[
		m=\rho V.
		\]
		
		\[
		m=
		7800\times6{,}54\times10^7
		\]
		
		\[
		\boxed{
			m_A\approx\SI{5.11e11}{kg}
		}
		\]
		
		\vspace{3mm}
		
		\textbf{39.}
		
		\[
		F_g=
		G\frac{m_Am_s}{r^2}
		\]
		
		\[
		F_g=
		\frac{
			(6{,}67430\times10^{-11})
			(5{,}11\times10^{11})
			(1000)
		}{
			(1000)^2
		}
		\]
		
		\[
		\boxed{
			F_g\approx\SI{3.41e-2}{N}
		}
		\]
		
		\vspace{3mm}
		
		\textbf{40.}
		
		Na Terra:
		
		\[
		P=mg
		\]
		
		\[
		P=1000\times9{,}8
		\]
		
		\[
		P=
		\SI{9.8e3}{N}.
		\]
		
		Comparando:
		
		\[
		\frac{9{,}8\times10^3}
		{3{,}41\times10^{-2}}
		\approx
		2{,}9\times10^5.
		\]
		
		Assim, o peso da sonda na Terra seria centenas de milhares de vezes
		superior à força gravítica exercida pelo asteroide na posição considerada.
		
		\vspace{3mm}
		
		\textbf{41.}
		
		Embora o asteroide pareça grande à escala humana, a sua massa é muito
		inferior à massa da Terra.
		
		Como
		
		\[
		F_g\propto M,
		\]
		
		a pequena massa do asteroide produz uma interação gravítica muito menos
		intensa.
		
	\end{solucaobox}
	
	% ============================================================
	\section*{Resolução 10 -- Quantidade de matéria}
	% ============================================================
	
	\begin{solucaobox}[Questões 42--46]
		
		\textbf{42.}
		
		\[
		m=nM
		\]
		
		\[
		m=
		2{,}00\times
		\SI{12.0}{g.mol^{-1}}
		\]
		
		\[
		m=\SI{24.0}{g}.
		\]
		
		Em quilogramas:
		
		\[
		\boxed{
			m=\SI{2.40e-2}{kg}
		}
		\]
		
		\vspace{3mm}
		
		\textbf{43.}
		
		\[
		N=nN_A
		\]
		
		\[
		N=
		2{,}00
		\times
		6{,}022\times10^{23}
		\]
		
		\[
		\boxed{
			N=
			1{,}2044\times10^{24}
			\text{ átomos}
		}
		\]
		
		\vspace{3mm}
		
		\textbf{44.}
		
		\[
		F_g=
		G\frac{m^2}{r^2}
		\]
		
		\[
		F_g=
		\frac{
			(6{,}67430\times10^{-11})
			(2{,}40\times10^{-2})^2
		}{
			(0{,}10)^2
		}
		\]
		
		\[
		\boxed{
			F_g\approx\SI{3.84e-12}{N}
		}
		\]
		
		\vspace{3mm}
		
		\textbf{45.}
		
		O valor obtido é extremamente pequeno.
		
		A constante gravítica $G$ é muito pequena e as massas de objetos
		laboratoriais são muito inferiores às massas astronómicas.
		
		Por isso, outras interações e forças presentes no laboratório dominam
		facilmente a interação gravítica.
		
		\vspace{3mm}
		
		\textbf{46.}
		
		Duplicar a quantidade de matéria duplica a massa de cada amostra.
		
		Assim:
		
		\[
		F'_g=
		G\frac{(2m)(2m)}{r^2}
		=
		4F_g.
		\]
		
		Logo, a força seria
		
		\[
		\boxed{4}
		\]
		
		vezes maior.
		
	\end{solucaobox}
	
	% ============================================================
	\section*{Resolução 11 -- Super-Terra}
	% ============================================================
	
	\begin{solucaobox}[Questões 47--51]
		
		\textbf{47.}
		
		Para a Terra:
		
		\[
		g_T=
		G\frac{M_T}{R_T^2}.
		\]
		
		Para o planeta:
		
		\[
		g_P=
		G\frac{5M_T}{(1{,}6R_T)^2}.
		\]
		
		Dividindo:
		
		\[
		\frac{g_P}{g_T}
		=
		\frac{
			G\dfrac{5M_T}{(1{,}6R_T)^2}
		}{
			G\dfrac{M_T}{R_T^2}
		}
		\]
		
		\[
		\boxed{
			\frac{g_P}{g_T}
			=
			\frac{5}{(1{,}6)^2}
		}
		\]
		
		\vspace{3mm}
		
		\textbf{48.}
		
		\[
		\frac{g_P}{g_T}
		=
		\frac{5}{2{,}56}
		\approx1{,}953.
		\]
		
		Logo,
		
		\[
		g_P
		\approx
		1{,}953\times9{,}82
		\]
		
		\[
		\boxed{
			g_P\approx\SI{19.2}{m.s^{-2}}
		}
		\]
		
		\vspace{3mm}
		
		\textbf{49.}
		
		\[
		P=mg
		\]
		
		\[
		P=
		70\times19{,}2
		\]
		
		\[
		\boxed{
			P\approx\SI{1.34e3}{N}
		}
		\]
		
		\vspace{3mm}
		
		\textbf{50.}
		
		Não.
		
		A massa da pessoa mantém-se em
		
		\[
		\boxed{
			\SI{70}{kg}
		}.
		\]
		
		\vspace{3mm}
		
		\textbf{51.}
		
		Não.
		
		O planeta possui uma massa cinco vezes superior, mas possui também um
		raio superior.
		
		Como
		
		\[
		g=
		\frac{GM}{R^2},
		\]
		
		o aumento do raio reduz a aceleração gravítica à superfície.
		
		Por isso, $g$ é apenas aproximadamente $1{,}95$ vezes maior que na Terra.
		
	\end{solucaobox}
	
	% ============================================================
	\section*{Resolução 12 -- Tratamento de dados}
	% ============================================================
	
	\begin{solucaobox}[Questões 52--56]
		
		Como
		
		\[
		m_1=m_2=\SI{1000}{kg},
		\]
		
		\[
		F_g=
		G\frac{10^6}{r^2}.
		\]
		
		Assim:
		
		\[
		Gm_1m_2
		=
		6{,}67430\times10^{-5}.
		\]
		
		\vspace{3mm}
		
		\textbf{52.}
		
		\[
		\begin{array}{c|c|c}
			r/\mathrm{m}
			&
			1/r^2\ /\mathrm{m^{-2}}
			&
			F_g/\mathrm{N}
			\\
			\hline
			1{,}0 & 1{,}000 & 6{,}67\times10^{-5}\\
			1{,}5 & 0{,}444 & 2{,}97\times10^{-5}\\
			2{,}0 & 0{,}250 & 1{,}67\times10^{-5}\\
			2{,}5 & 0{,}160 & 1{,}07\times10^{-5}\\
			3{,}0 & 0{,}111 & 7{,}42\times10^{-6}
		\end{array}
		\]
		
		\vspace{3mm}
		
		\textbf{53.}
		
		O gráfico $F_g=f(r)$ apresenta uma curva decrescente semelhante à
		representada anteriormente na ficha.
		
		\vspace{3mm}
		
		\textbf{54.}
		
		Não.
		
		A força não varia linearmente com a distância:
		
		\[
		F_g\propto\frac1{r^2}.
		\]
		
		\vspace{3mm}
		
		\textbf{55.}
		
		Definindo
		
		\[
		x=\frac1{r^2},
		\]
		
		tem-se
		
		\[
		F_g=
		Gm_1m_2x.
		\]
		
		Portanto, o gráfico
		
		\[
		F_g=f\left(\frac1{r^2}\right)
		\]
		
		é uma reta que passa pela origem.
		
		\vspace{3mm}
		
		\textbf{56.}
		
		O declive é:
		
		\[
		k=Gm_1m_2.
		\]
		
		Logo,
		
		\[
		k=
		(6{,}67430\times10^{-11})
		(1000)(1000)
		\]
		
		\[
		\boxed{
			k=
			6{,}67430\times10^{-5}
		}
		\]
		
		em unidades compatíveis com o gráfico.
		
	\end{solucaobox}
	
	% ============================================================
	\section*{Resolução 13 -- Órbita circular}
	% ============================================================
	
	\begin{solucaobox}[Questões 57--61]
		
		\textbf{57.}
		
		Para uma órbita circular:
		
		\[
		F_g=F_R.
		\]
		
		Então,
		
		\[
		G\frac{M_Tm}{r^2}
		=
		m\frac{v^2}{r}.
		\]
		
		Cancelando $m$:
		
		\[
		G\frac{M_T}{r^2}
		=
		\frac{v^2}{r}.
		\]
		
		Multiplicando por $r$:
		
		\[
		\frac{GM_T}{r}
		=
		v^2.
		\]
		
		Logo,
		
		\[
		\boxed{
			v=
			\sqrt{\frac{GM_T}{r}}
		}
		\]
		
		\vspace{3mm}
		
		\textbf{58.}
		
		Com
		
		\[
		r=\SI{6.771e6}{m},
		\]
		
		\[
		v=
		\sqrt{
			\frac{
				(6{,}67430\times10^{-11})
				(5{,}972\times10^{24})
			}{
				6{,}771\times10^6
			}
		}.
		\]
		
		\[
		\boxed{
			v\approx\SI{7.67e3}{m.s^{-1}}
		}
		\]
		
		ou
		
		\[
		\boxed{
			v\approx\SI{7.67}{km.s^{-1}}
		}
		\]
		
		\vspace{3mm}
		
		\textbf{59.}
		
		\[
		T=
		\frac{2\pi r}{v}
		\]
		
		\[
		T=
		\frac{
			2\pi(6{,}771\times10^6)
		}{
			7{,}672\times10^3
		}
		\]
		
		\[
		T\approx\SI{5.54e3}{s}.
		\]
		
		Em minutos:
		
		\[
		T=
		\frac{5545}{60}
		\]
		
		\[
		\boxed{
			T\approx\SI{92.4}{min}
		}
		\]
		
		\vspace{3mm}
		
		\textbf{60.}
		
		\[
		E_c=
		\frac12mv^2
		\]
		
		\[
		E_c=
		\frac12
		(1200)
		(7{,}672\times10^3)^2
		\]
		
		\[
		\boxed{
			E_c\approx\SI{3.53e10}{J}
		}
		\]
		
		\vspace{3mm}
		
		\textbf{61.}
		
		A massa do satélite aparece tanto na força gravítica como na expressão
		da força resultante centrípeta.
		
		Ao igualar:
		
		\[
		G\frac{M_Tm}{r^2}
		=
		m\frac{v^2}{r},
		\]
		
		a massa $m$ cancela-se.
		
		Por isso, para uma mesma órbita, a velocidade orbital não depende da
		massa do satélite.
		
	\end{solucaobox}
	
	% ============================================================
	\section*{Resolução 14 -- Escolha múltipla}
	% ============================================================
	
	\begin{solucaobox}[Questões 62--67]
		
		\textbf{62.} B
		
		\[
		F_g\propto m_1m_2.
		\]
		
		Duplicar uma massa duplica a força.
		
		\vspace{2mm}
		
		\textbf{63.} B
		
		\[
		(2m_1)(2m_2)=4m_1m_2.
		\]
		
		\vspace{2mm}
		
		\textbf{64.} D
		
		\[
		F'_g=
		\frac{F_g}{3^2}
		=
		\frac{F_g}{9}.
		\]
		
		\vspace{2mm}
		
		\textbf{65.} C
		
		As forças de interação têm a mesma intensidade e sentidos opostos.
		
		\vspace{2mm}
		
		\textbf{66.} C
		
		A massa do astronauta mantém-se.
		
		\vspace{2mm}
		
		\textbf{67.} C
		
		\[
		\boxed{
			g=
			\frac{GM}{r^2}
		}
		\]
		
	\end{solucaobox}
	
	% ============================================================
	\section*{Resolução 15 -- Problema final}
	% ============================================================
	
	\begin{solucaobox}[Questões 68--74]
		
		Dados:
		
		\[
		M=\SI{8.0e23}{kg}
		\]
		
		\[
		R=\SI{4.0e6}{m}
		\]
		
		\[
		m=\SI{500}{kg}.
		\]
		
		\vspace{3mm}
		
		\textbf{68.}
		
		\[
		g=
		G\frac{M}{R^2}
		\]
		
		\[
		g=
		\frac{
			(6{,}67430\times10^{-11})
			(8{,}0\times10^{23})
		}{
			(4{,}0\times10^6)^2
		}
		\]
		
		\[
		g=
		\frac{
			5{,}33944\times10^{13}
		}{
			1{,}6\times10^{13}
		}
		\]
		
		\[
		\boxed{
			g\approx\SI{3.34}{m.s^{-2}}
		}
		\]
		
		\vspace{3mm}
		
		\textbf{69.}
		
		\[
		P=mg
		\]
		
		\[
		P=
		500\times3{,}34
		\]
		
		\[
		\boxed{
			P\approx\SI{1.67e3}{N}
		}
		\]
		
		\vspace{3mm}
		
		\textbf{70.}
		
		Se
		
		\[
		h=R,
		\]
		
		então
		
		\[
		r=R+h=2R.
		\]
		
		Logo,
		
		\[
		\boxed{
			r=\SI{8.0e6}{m}
		}
		\]
		
		\vspace{3mm}
		
		\textbf{71.}
		
		Como a distância ao centro duplicou:
		
		\[
		g'=
		\frac{g}{2^2}
		=
		\frac{3{,}34}{4}.
		\]
		
		\[
		\boxed{
			g'\approx\SI{0.834}{m.s^{-2}}
		}
		\]
		
		\vspace{3mm}
		
		\textbf{72.}
		
		\[
		F'_g=mg'
		\]
		
		\[
		F'_g=
		500\times0{,}834
		\]
		
		\[
		\boxed{
			F'_g\approx\SI{417}{N}
		}
		\]
		
		\vspace{3mm}
		
		\textbf{73.}
		
		\[
		\frac{
			F_{\text{altitude}}
		}{
			F_{\text{superfície}}
		}
		=
		\frac{417}{1670}
		\]
		
		\[
		\boxed{
			\frac{
				F_{\text{altitude}}
			}{
				F_{\text{superfície}}
			}
			\approx0{,}25
		}
		\]
		
		\vspace{3mm}
		
		\textbf{74.}
		
		Ao subir para uma altitude igual ao raio do planeta, a distância ao centro
		passa de $R$ para $2R$.
		
		Como
		
		\[
		F_g\propto\frac1{r^2},
		\]
		
		então
		
		\[
		\frac{F'_g}{F_g}
		=
		\frac{R^2}{(2R)^2}
		=
		\frac14.
		\]
		
		Assim, a força gravítica reduz-se para um quarto do valor à superfície.
		
	\end{solucaobox}
	
	% ============================================================
	% FECHO
	% ============================================================
	
	\clearpage
	
	\begin{center}
		
		\begin{tikzpicture}
			
			\node[
			rounded corners=3mm,
			draw=eufblue,
			line width=1pt,
			fill=euflight,
			inner xsep=10mm,
			inner ysep=5mm
			]
			{
				
				\begin{minipage}{0.84\textwidth}
					
					\centering
					
					{\Large\bfseries\textcolor{eufblue}{EUFISICA}}
					
					\vspace{2mm}
					
					\textbf{Lei da Gravitação Universal}
					
					\vspace{2mm}
					
					\small
					
					Massa $\bullet$
					Distância $\bullet$
					Força $\bullet$
					Campo gravítico $\bullet$
					Órbitas $\bullet$
					Exploração espacial
					
					\vspace{2mm}
					
					\textit{
						A mesma lei física que explica a queda de um corpo
						permite compreender o movimento de luas, planetas e satélites.
					}
					
					\vspace{2mm}
					
					\textbf{Prof. José Gonçalves}
					
				\end{minipage}
				
			};
			
		\end{tikzpicture}
		
	\end{center}
	
	\fi
	
\end{document}